% ====== Tipo de archivo ======
%\documentclass[journal]{IEEEtran}
\documentclass{article}

% ====== Paquetes básicos ======
\usepackage[utf8]{inputenc}    % Codificación UTF-8
\usepackage[T1]{fontenc}       % Acentos correctos
\usepackage[spanish]{babel}    % Traducción al español
\usepackage{csquotes}          % Recomendado para biblatex con babel
\usepackage{hyperref}          % Hipervínculos
\usepackage{amsmath,amssymb}   % Matemáticas
\usepackage{graphicx}          % Imágenes
\usepackage{tikz}              % Diagramas vectoriales
\usetikzlibrary{arrows.meta,positioning}
\usepackage{caption}           % Mejor control de captions
\usepackage{geometry}          % Márgenes
\usepackage{fancyhdr}          % Encabezados y pies
\usepackage{parskip}           % Mejor separación entre párrafos

% ====== Referencias ======
\usepackage[
    backend=biber,
    style=numeric-comp,     % cambie a APA de ser necesario
    sorting=none
]{biblatex} 
\bibliography{references}

% agrege `\printbibliography{}` para las referencias
% cite con `\cite{}`
% agregue las referencias a un `references.bib`

% ====== Configuración ======
\geometry{margin=2.5cm}
\pagestyle{fancy}
\fancyhf{}
\rhead{\thepage}
\setlength{\headheight}{13pt}
\lhead{\leftmark}

% ====== Comando personalizado para imágenes ======
% Uso: \imagen{ruta}{Texto del pie de figura}
\newcommand{\imagen}[2]{
    \begin{figure}[h!]
        \centering
        \includegraphics[width=200px]{#1}
        \caption{#2}
    \end{figure}
}

% ====== Datos del artículo ======
\title{Portada}
\author{Lester Iván Hernández Hernandez}


\usepackage[spanish]{babel}
\usepackage{datetime}
\newdateformat{mydateformat}{\THEDAY/\THEMONTH/\THEYEAR}

\newcommand{\titlename}{Asignación 3er Parcial}
\newcommand{\signature}{Clase (QQ-100)}
\newcommand{\seccion}{1100}
\newcommand{\docente}{Alex Conrado Mejía Valladares}
%Individual
\newcommand{\alumno}{Lester Iván Hernández Hernandez}
\newcommand{\cuenta}{20242000637}

%Grupo
%\newcommand{\grupo}{2}
%\newcommand{\integrantes}{}

\newcommand{\portada}{
    \begin{titlepage}
        \centering
        \vspace*{-1cm}
        \includegraphics[width=4.5cm]{Logo-unah.jpg}\\
        \vspace*{0.5cm}
        {\large\bfseries UNIVERSIDAD NACIONAL AUTÓNOMA DE HONDURAS
        \\UNAH-VS\par}
        \vspace{2cm}
        {\large \titlename\par}
        \vspace{1cm}
        {\large Asignatura: \signature\par}
        {\large Sección: \seccion\par}
        {\large Docente: \docente\par}
        \vspace{1cm}
        %Individual
        {\large Alumno: \alumno\par}
        {\large N. de cuenta: \cuenta\par}
        %Grupo
        %{\large Grupo: \grupo\par}
        %{\large Integrantes: \integrantes\par}
        \vfill
        {\large \today\par}
        \vspace*{1in}
    \end{titlepage}
}


\begin{document}
    \portada{}
    \newpage
    
    % =========================================================
    %  PROCESOS DE PRODUCCIÓN DE GASES
    % =========================================================

    \section{Producción de Hidrógeno (H\textsubscript{2})}

    El hidrógeno es el elemento más abundante del universo, pero en la Tierra
    prácticamente no se encuentra en estado libre. Para obtenerlo es necesario
    separarlo de los compuestos que lo contienen, como el agua o el gas natural.
    El método industrial dominante es el \textbf{reformado de vapor de metano}
    (\textit{Steam Methane Reforming}, SMR), responsable de más del 90\% de la
    producción mundial.

    El proceso comienza con una etapa de \textbf{desulfuración}: el gas natural
    (cuyo componente principal es el metano, CH\textsubscript{4}) se hace pasar
    por un lecho de material absorbente que elimina los compuestos de azufre. Este
    paso es indispensable porque el azufre envenena los catalizadores metálicos
    usados en las etapas siguientes, reduciendo drásticamente su vida útil.

    A continuación, el gas limpio se mezcla con vapor de agua sobrecalentado y se
    introduce en un reactor tubular a temperaturas de entre 700 y 1\,000\,°C, en
    presencia de un catalizador de níquel. Bajo estas condiciones, el metano
    reacciona con el vapor produciendo hidrógeno y monóxido de carbono (CO):
    \begin{equation}
        \text{CH}_4 + \text{H}_2\text{O} \;\rightarrow\; \text{CO} + 3\,\text{H}_2
    \end{equation}

    El CO resultante no se descarta: reacciona con vapor adicional en dos reactores
    en serie —uno a alta temperatura y otro a baja temperatura— en la llamada
    \textbf{reacción de desplazamiento de agua} (\textit{Water-Gas Shift}), que
    genera más hidrógeno y convierte el CO en dióxido de carbono:
    \begin{equation}
        \text{CO} + \text{H}_2\text{O} \;\rightarrow\; \text{CO}_2 + \text{H}_2
    \end{equation}

    Finalmente, la corriente gaseosa pasa por una unidad de \textbf{adsorción por
    variación de presión} (PSA), donde las impurezas quedan retenidas en un
    material sólido poroso y el hidrógeno sale con una pureza superior al 99.9\%.
    El producto se emplea en la fabricación de fertilizantes, el refinado de
    petróleo y, cada vez más, como combustible limpio en celdas de combustible.

    % ---------------------------------------------------------

    \section{Producción de Amoniaco (NH\textsubscript{3})}

    El amoniaco es una molécula compuesta por nitrógeno e hidrógeno que tiene un
    olor fuerte y característico. Es la base de la mayoría de los fertilizantes
    nitrogenados y, por tanto, uno de los productos químicos más importantes para
    la agricultura moderna. Se sintetiza industrialmente mediante el
    \textbf{proceso Haber-Bosch}, desarrollado a principios del siglo XX y
    considerado uno de los avances tecnológicos más influyentes de la historia, ya
    que permitió producir alimentos para una población mundial en rápido crecimiento.

    Las materias primas son el nitrógeno molecular (N\textsubscript{2}), obtenido
    directamente del aire por destilación criogénica, y el hidrógeno
    (H\textsubscript{2}), proveniente generalmente del proceso SMR descrito
    anteriormente. Ambos gases se purifican con cuidado para eliminar trazas de
    oxígeno, agua y CO que desactivarían el catalizador, y se mezclan en una
    proporción molar de 1:3.

    La mezcla se comprime hasta presiones de entre 150 y 300 atmósferas y se
    hace circular a través de un reactor que contiene un catalizador de hierro
    promovido con pequeñas cantidades de potasio y aluminio, operando a una
    temperatura de 400--500\,°C:
    \begin{equation}
        \text{N}_2 + 3\,\text{H}_2 \;\rightleftharpoons\; 2\,\text{NH}_3
    \end{equation}

    Esta combinación de alta presión y temperatura moderada es un compromiso
    deliberado: la presión alta favorece la formación de amoniaco, mientras que
    la temperatura no puede ser demasiado elevada porque, a mayor temperatura,
    el equilibrio se desplaza hacia los reactivos. El catalizador de hierro
    permite alcanzar velocidades de reacción aceptables sin necesidad de
    temperaturas excesivas. Aun así, solo entre el 15 y el 25\% del gas se
    convierte en amoniaco por pasada, de modo que la mezcla se enfría para
    licuar el amoniaco producido y los gases no reaccionados se recirculan al
    reactor. El amoniaco líquido resultante se almacena en tanques refrigerados
    a $-33$\,°C. Su principal destino es la producción de urea y otros
    fertilizantes, aunque también se usa como refrigerante industrial y en la
    síntesis de explosivos y plásticos.

    % ---------------------------------------------------------

    \section{Producción de Cloro (Cl\textsubscript{2})}

    El cloro es un gas de color amarillo-verdoso con un olor fuerte e irritante.
    Es fundamental para desinfectar el agua potable, fabricar plásticos como el
    PVC y producir una amplia variedad de productos farmacéuticos y agroquímicos.
    El proceso industrial dominante para obtenerlo es la \textbf{electrólisis de
    la salmuera}, conocido como \textbf{proceso cloroalcalino}.

    La materia prima es una solución saturada de cloruro de sodio (NaCl) en agua,
    llamada \textit{salmuera}. Antes de ingresar a la celda electrolítica, la
    salmuera se purifica exhaustivamente para eliminar iones de calcio, magnesio
    y sulfato, que deteriorarían la membrana selectiva del reactor.

    La celda está dividida en dos compartimentos por una \textbf{membrana de
    Nafion\textregistered{}}, un polímero que permite el paso de iones
    Na\textsuperscript{+} de un lado al otro, pero no el de cloro ni hidróxido.
    Al aplicar corriente eléctrica, ocurren dos reacciones simultáneas en los
    electrodos:
    \begin{equation}
        \text{Ánodo (+): } 2\,\text{Cl}^- \;\rightarrow\; \text{Cl}_2 + 2\,e^-
    \end{equation}
    \begin{equation}
        \text{Cátodo (-): } 2\,\text{H}_2\text{O} + 2\,e^- \;\rightarrow\;
        \text{H}_2 + 2\,\text{OH}^-
    \end{equation}

    En el ánodo, los iones cloruro son oxidados y se libera gas cloro. En el
    cátodo, el agua se reduce produciendo gas hidrógeno e iones hidróxido. Los
    iones Na\textsuperscript{+} que cruzan la membrana se combinan con los
    OH\textsuperscript{-} del cátodo para formar hidróxido de sodio (NaOH) en
    solución. De esta manera, un solo proceso genera simultáneamente tres
    productos de valor comercial: cloro, sosa cáustica e hidrógeno.

    El gas cloro que sale del ánodo se enfría para condensar el agua que arrastra,
    se seca con ácido sulfúrico concentrado y finalmente se comprime y licua para
    facilitar su transporte. El hidrógeno se usa como combustible dentro de la
    misma planta o se vende, y el NaOH en solución se concentra por evaporación.
    Este proceso consume grandes cantidades de electricidad, por lo que las plantas
    cloroalcalinas generalmente se ubican cerca de fuentes de energía económica.



\end{document} 